Sensitivity to the assumed washing threshold, ground market. Rectangular band, h = 60 mm, f = 0.1, initial added margin 0 mm; all other assumptions fixed. Columns: Neuchâtel, Almería and Doha. a–c, additional annual energy loss due to the band without manual cleaning: loss with band minus loss without band, as a percentage of identical clean production. d–f, net present value of an additional 10 mm margin over 25 years at a 5\% discount rate, in EUR/kWp, including the extra-glass investment; costs with and without band are independently optimised at each margin, as in the existing economic analysis. g–i, annual net gain from manual cleaning for total soiling (uniform deposit plus band), as a percentage of clean revenue: total cost with no manual cleaning minus the minimum total cost of the evaluated manual policies, including cleaning charges. This selection considers economic-trigger and fixed-interval manual policies (14, 30, 60, 90, 120, 180, 270 and 365 days), excluding both no manual cleaning and the robot; negative gains are retained. Colours identify modules; solid lines/circles denote portrait and dashed lines/triangles landscape. The dotted vertical line marks $R_H$ = 15 mm/event, an illustrative reference specified for the main analysis, not a threshold chosen by optimising these curves. Horizontal lines mark zero economic benefit. Panels a, g use a separate Neuchâtel scale; b–c and h–i share scales, and d–f use a common scale. Curves show individual configurations; no uncertainty bars or confidence intervals are used. Complete removal of uniform and band deposits is assumed when event precipitation $P_e$ $\geq$ $R_H$. Events are separated by three consecutive dry hours. These washing thresholds are not calibrated for bottom-edge bands. The annual normalisation and available policy sets follow the saved simulations.
